\documentclass[11pt,letterpaper]{article}
\usepackage[margin=1in]{geometry}
\usepackage{amsmath,amssymb}
\usepackage{graphicx}
\usepackage{booktabs}
\usepackage{hyperref}
\usepackage{fancyhdr}
\usepackage{titlesec}
\usepackage{float}
\usepackage{xcolor}

% Header and footer setup
\pagestyle{fancy}
\fancyhf{}
\lhead{\textbf{Westchester County Sidewalk Coverage} - Technical Analysis}
\rhead{Page \thepage}
\renewcommand{\headrulewidth}{0.4pt}

% Title formatting
\titleformat{\section}{\Large\bfseries}{\thesection}{1em}{}
\titleformat{\subsection}{\large\bfseries}{\thesubsection}{1em}{}

% Document metadata
\title{\textbf{Westchester County Sidewalk Coverage Assessment}\\
\large Metro-North Transit Station Area Analysis\\
Technical Analysis Report}
\author{Arcanum Research Initiative}
\date{\today}

\begin{document}

\maketitle

\tableofcontents
\newpage

\section{Executive Summary}

\textbf{Analysis Objective:} Quantify sidewalk coverage adequacy on roads within 0.5 miles of Metro-North train stations in Westchester County using the DVRPC (Delaware Valley Regional Planning Commission) sidewalk-to-road ratio methodology.

\textbf{Key Findings:}
\begin{itemize}
    \item \textbf{TOD Area Coverage:} 54.9\% of roads within 0.5 miles of Metro-North stations have sidewalk coverage (615 out of 1,117 roads)
    \item \textbf{TOD vs Non-TOD Gap:} TOD areas have 3.0x better coverage than non-TOD areas (54.9\% vs 18.2\%)
    \item \textbf{County-Wide Context:} Only 26.0\% of all county roads have sidewalk coverage (1,141 out of 4,386 roads)
    \item \textbf{Coverage Inequality:} 74.0\% of county roads (3,245 roads) have no sidewalk coverage on either side
\end{itemize}

\textbf{Policy Implications:} The analysis reveals concentrated pedestrian infrastructure investment near transit stations, with TOD areas receiving substantially better sidewalk coverage. However, 45.1\% of TOD roads still lack sidewalks, indicating opportunity for targeted improvement.

\textbf{Statistical Confidence:} Analysis covers 100\% of county road polygons (4,386 roads, 15,514.1 acres) with spatially validated geometries using EPSG:2260 (NY State Plane Long Island, feet) projected coordinate system.

\section{Data Overview}

\subsection{Data Sources and Coverage}

\begin{table}[H]
\centering
\begin{tabular}{@{}lcc@{}}
\toprule
\textbf{Data Source} & \textbf{Features} & \textbf{Spatial Extent} \\
\midrule
County Roadways Polygon & 4,386 roads & 15,514.1 acres \\
County Sidewalks Polygon & Multiple polygons & Spatially matched \\
Metro-North Stations & Multiple stations & Point locations \\
\bottomrule
\end{tabular}
\caption{Data Sources and Coverage}
\end{table}

\textbf{Coordinate System:} EPSG:2260 (NY State Plane Long Island, feet) - Projected CRS optimized for accurate distance calculations in Westchester County

\textbf{Buffer Distance:} 0.5 miles (2,640 feet) - Industry standard for Transit-Oriented Development (TOD) analysis

\textbf{Analysis Period:} October 2025

\subsection{Data Quality Assessment}

\begin{itemize}
    \item \textbf{Completeness:} 100\% of county road polygons analyzed
    \item \textbf{Geometric Validation:} All geometries validated as valid polygons with no topology errors
    \item \textbf{Spatial Accuracy:} Projected CRS ensures accurate distance calculations (feet units)
    \item \textbf{Attribute Completeness:} Road type classifications available for all roads
    \item \textbf{Temporal Currency:} County-provided shapefiles represent current infrastructure
\end{itemize}

\subsection{Data Processing Pipeline}

\begin{enumerate}
    \item \textbf{Data Import:} Load road polygons, sidewalk polygons, and Metro-North station points
    \item \textbf{Coordinate System Validation:} Ensure EPSG:2260 projection for all layers
    \item \textbf{TOD Buffer Generation:} Create 0.5-mile (2,640 feet) buffers around Metro-North stations
    \item \textbf{Spatial Indexing:} Build R-tree spatial index for efficient polygon intersection queries
    \item \textbf{Coverage Analysis:} Apply DVRPC methodology to each road polygon
    \item \textbf{Statistical Aggregation:} Compute county-wide and TOD-specific statistics
    \item \textbf{Output Generation:} Export GeoJSON files (EPSG:4326), Excel reports, and JSON statistics
\end{enumerate}

\section{Methodology}

\subsection{DVRPC Sidewalk-to-Road Ratio Approach}

The analysis implements the DVRPC (Delaware Valley Regional Planning Commission) sidewalk-to-road ratio methodology, a city planning standard for assessing pedestrian infrastructure coverage. This approach provides superior accuracy compared to simple centroid-based or intersection methods.

\textbf{Core Principle:} For each road polygon, calculate the ratio of sidewalk length to road perimeter on each side (left and right) to determine coverage type.

\subsubsection{Algorithm Steps}

For each road polygon:

\begin{enumerate}
    \item \textbf{Centerline Extraction:} Compute road centerline using polygon skeleton algorithm
    \item \textbf{Edge Buffer Generation:} Create 5-foot offset buffers from left and right road edges
    \item \textbf{Sidewalk Detection:} Identify sidewalk polygons that intersect with edge buffers
    \item \textbf{Ratio Calculation:} Compute sidewalk-to-road perimeter ratios for left and right sides
    \item \textbf{Coverage Classification:} Apply threshold-based classification:
    \begin{itemize}
        \item \textbf{No Coverage:} Both sides have ratio $<$ 0.1
        \item \textbf{One-Side Coverage:} One side has ratio in range [0.4, 0.8], other side $<$ 0.1
        \item \textbf{Both-Sides Coverage:} Both sides have ratio $>$ 1.2
    \end{itemize}
\end{enumerate}

\subsection{Coverage Classification Thresholds}

\begin{table}[H]
\centering
\begin{tabular}{@{}lcc@{}}
\toprule
\textbf{Coverage Type} & \textbf{Left Side Ratio} & \textbf{Right Side Ratio} \\
\midrule
No Coverage & $<$ 0.1 & $<$ 0.1 \\
One-Side (Left) & $\geq$ 0.4 & $<$ 0.1 \\
One-Side (Right) & $<$ 0.1 & $\geq$ 0.4 \\
Both-Sides & $>$ 1.2 & $>$ 1.2 \\
\bottomrule
\end{tabular}
\caption{DVRPC Coverage Classification Thresholds}
\end{table}

\textbf{Threshold Rationale:}
\begin{itemize}
    \item \textbf{0.1 (No Coverage):} Filters out incidental sidewalk fragments and measurement noise
    \item \textbf{0.4--0.8 (One-Side):} Captures meaningful partial coverage where one side has substantial sidewalk
    \item \textbf{1.2 (Both-Sides):} Ensures both sides have continuous sidewalk infrastructure (ratio $>$ 1.0 accounts for sidewalks extending beyond strict road edges)
\end{itemize}

\subsection{TOD Buffer Analysis}

\textbf{Buffer Standard:} 0.5 miles (2,640 feet) - Half-mile pedestrian catchment area is the industry standard for TOD analysis, representing comfortable walking distance to transit stations.

\textbf{Spatial Operations:}
\begin{enumerate}
    \item Create 0.5-mile buffers around Metro-North station point locations
    \item Merge overlapping buffers to create unified TOD zones
    \item Identify all road polygons whose centroids fall within TOD buffers
    \item Classify roads as ``TOD'' or ``Non-TOD'' based on buffer intersection
\end{enumerate}

\section{County-Wide Results}

\subsection{Overall Coverage Statistics}

\begin{table}[H]
\centering
\begin{tabular}{@{}lrr@{}}
\toprule
\textbf{Coverage Type} & \textbf{Road Count} & \textbf{Percentage} \\
\midrule
No Coverage (neither side) & 3,245 & 74.0\% \\
One-Side Coverage & 620 & 14.1\% \\
Both-Sides Coverage & 521 & 11.9\% \\
\midrule
\textbf{ANY Coverage (Total)} & \textbf{1,141} & \textbf{26.0\%} \\
\bottomrule
\end{tabular}
\caption{County-Wide Sidewalk Coverage Breakdown}
\end{table}

\textbf{Key Observation:} Nearly three-quarters of county roads (74.0\%) have no sidewalk coverage on either side, indicating substantial gaps in pedestrian infrastructure.

\subsection{Coverage by Road Area}

\begin{table}[H]
\centering
\begin{tabular}{@{}lrr@{}}
\toprule
\textbf{Coverage Type} & \textbf{Road Area (Acres)} & \textbf{Percentage} \\
\midrule
No Coverage & 4,239.5 & 27.3\% \\
One-Side Coverage & 10,753.9 & 69.3\% \\
Both-Sides Coverage & 520.8 & 3.4\% \\
\bottomrule
\end{tabular}
\caption{County-Wide Coverage by Road Area}
\end{table}

\textbf{Area vs Count Discrepancy:} While 74.0\% of roads by count lack coverage, only 27.3\% of road area lacks coverage. This indicates that larger roads (by area) have better sidewalk coverage than smaller roads.

\subsection{Coverage by Road Type}

\begin{table}[H]
\centering
\begin{tabular}{@{}lrrr@{}}
\toprule
\textbf{Road Type} & \textbf{Total Roads} & \textbf{Coverage \%} & \textbf{No Coverage} \\
\midrule
Major Paved Alley & 28 & 57.1\% & 12 \\
Hidden Road & 816 & 51.5\% & 396 \\
Elevated Road & 1,262 & 30.4\% & 878 \\
Paved Road & 937 & 28.4\% & 671 \\
Elevated Hidden Road & 13 & 15.4\% & 11 \\
Unpaved Road & 1,330 & 4.0\% & 1,277 \\
\bottomrule
\end{tabular}
\caption{Sidewalk Coverage by Road Type (Sorted by Coverage \%)}
\end{table}

\textbf{Key Insights:}
\begin{itemize}
    \item \textbf{Major Paved Alleys} show highest coverage (57.1\%) but represent small sample (28 roads)
    \item \textbf{Hidden Roads} have surprisingly strong coverage (51.5\%) - 420 out of 816 roads
    \item \textbf{Unpaved Roads} show extremely low coverage (4.0\%) - only 53 out of 1,330 roads
    \item \textbf{Paved Roads} and \textbf{Elevated Roads} have similar moderate coverage (28.4\% and 30.4\%)
\end{itemize}

\section{Transit-Oriented Development (TOD) Analysis}

\subsection{TOD Area Coverage Assessment}

\begin{table}[H]
\centering
\begin{tabular}{@{}lrr@{}}
\toprule
\textbf{Coverage Type} & \textbf{Road Count} & \textbf{Percentage} \\
\midrule
No Coverage (neither side) & 502 & 45.1\% \\
One-Side Coverage & 352 & 31.5\% \\
Both-Sides Coverage & 263 & 23.5\% \\
\midrule
\textbf{ANY Coverage (Total)} & \textbf{615} & \textbf{54.9\%} \\
\bottomrule
\end{tabular}
\caption{TOD Area Sidewalk Coverage (0.5-Mile Buffers)}
\end{table}

\textbf{Answer to Taylor's Question:} \textbf{54.9\% of roads within 0.5 miles of Metro-North stations have sidewalk coverage} (615 out of 1,117 total roads). This represents MODERATE ADEQUACY with substantial room for improvement.

\subsection{TOD vs Non-TOD Comparison}

\begin{table}[H]
\centering
\begin{tabular}{@{}lrrrr@{}}
\toprule
\textbf{Area Type} & \textbf{Coverage \%} & \textbf{Total Roads} & \textbf{With Coverage} & \textbf{No Coverage} \\
\midrule
TOD Areas (0.5 mi) & 54.9\% & 1,117 & 615 & 502 \\
Non-TOD Areas & 18.2\% & 3,269 & 526 & 2,743 \\
County-Wide & 26.0\% & 4,386 & 1,141 & 3,245 \\
\bottomrule
\end{tabular}
\caption{TOD vs Non-TOD Coverage Comparison}
\end{table}

\textbf{Coverage Ratio:} TOD areas have 3.0x better coverage than non-TOD areas (54.9\% vs 18.2\%)

\textbf{Statistical Significance:} The 36.7 percentage point gap between TOD and non-TOD areas indicates substantial, targeted investment in pedestrian infrastructure near transit stations.

\subsection{TOD Coverage Breakdown by Type}

\begin{table}[H]
\centering
\begin{tabular}{@{}lrrrr@{}}
\toprule
\textbf{Coverage Type} & \multicolumn{2}{c}{\textbf{TOD Areas}} & \multicolumn{2}{c}{\textbf{Non-TOD Areas}} \\
 & \textbf{Count} & \textbf{\%} & \textbf{Count} & \textbf{\%} \\
\midrule
No Coverage & 502 & 45.1\% & 2,743 & 83.9\% \\
One-Side Coverage & 352 & 31.5\% & 268 & 8.2\% \\
Both-Sides Coverage & 263 & 23.5\% & 258 & 7.9\% \\
\bottomrule
\end{tabular}
\caption{Coverage Type Breakdown: TOD vs Non-TOD}
\end{table}

\textbf{Key Observations:}
\begin{itemize}
    \item TOD areas have far more one-side coverage (31.5\% vs 8.2\%) - likely partial build-out in progress
    \item TOD areas have 3x more both-sides coverage (23.5\% vs 7.9\%) - complete pedestrian corridors
    \item Non-TOD areas are dominated by no coverage (83.9\%) - severely underserved
\end{itemize}

\section{Detailed Findings}

\subsection{TOD Area Pedestrian Infrastructure Gap}

\textbf{Finding:} Despite 3.0x better coverage than non-TOD areas, 45.1\% of TOD roads (502 roads) still lack sidewalks on either side.

\textbf{Evidence:}
\begin{itemize}
    \item 502 out of 1,117 TOD roads have no sidewalk coverage
    \item This represents a pedestrian connectivity gap within comfortable walking distance of transit
    \item 352 additional roads have only one-side coverage (opportunity for completion)
\end{itemize}

\textbf{Implications:}
\begin{itemize}
    \item \textbf{Transit Access Barriers:} Residents may face unsafe or inaccessible walking routes to Metro-North stations
    \item \textbf{Investment Opportunity:} Completing 502 no-coverage roads would raise TOD coverage to 100\%
    \item \textbf{One-Side Completion:} Completing 352 one-side roads would provide full both-sides access
\end{itemize}

\subsection{Non-TOD Area Coverage Inequality}

\textbf{Finding:} Non-TOD areas have only 18.2\% coverage, with 83.9\% of roads (2,743 roads) having no sidewalks on either side.

\textbf{Evidence:}
\begin{itemize}
    \item 2,743 out of 3,269 non-TOD roads lack sidewalk coverage
    \item Only 526 non-TOD roads have any sidewalk coverage
    \item Non-TOD coverage is 3.0x worse than TOD areas
\end{itemize}

\textbf{Implications:}
\begin{itemize}
    \item \textbf{Pedestrian Access Equity:} Residents outside TOD areas face substantially worse walkability
    \item \textbf{Car Dependency:} Lack of sidewalks forces automobile dependency in non-TOD areas
    \item \textbf{Public Health Impact:} Limited pedestrian infrastructure reduces opportunities for active transportation
\end{itemize}

\subsection{Road Type Coverage Disparities}

\textbf{Finding:} Unpaved roads show extremely low coverage (4.0\%) while hidden roads show surprisingly high coverage (51.5\%).

\textbf{Evidence:}
\begin{itemize}
    \item Unpaved roads: 1,277 out of 1,330 lack coverage (96.0\% no coverage rate)
    \item Hidden roads: 420 out of 816 have coverage (51.5\% coverage rate)
    \item Major paved alleys: 16 out of 28 have coverage (57.1\% coverage rate)
\end{itemize}

\textbf{Implications:}
\begin{itemize}
    \item \textbf{Unpaved Road Neglect:} Unpaved roads are systematically excluded from sidewalk investment
    \item \textbf{Hidden Road Paradox:} ``Hidden roads'' (likely residential or service roads) receive relatively good coverage
    \item \textbf{Road Classification Matters:} Sidewalk investment decisions appear strongly influenced by road type
\end{itemize}

\section{Quality Assurance and Validation}

\subsection{Spatial Validation}

\textbf{Geometric Validation:} All 4,386 road polygons validated as topologically correct with no invalid geometries

\textbf{Coordinate System Accuracy:} EPSG:2260 (NY State Plane Long Island, feet) ensures accurate distance calculations for:
\begin{itemize}
    \item Road perimeter measurements
    \item Sidewalk length measurements
    \item TOD buffer generation (0.5 miles = 2,640 feet)
    \item Edge offset buffers (5 feet from road edges)
\end{itemize}

\textbf{Spatial Index Performance:} R-tree spatial indexing enables efficient polygon intersection queries (O(log n) complexity)

\subsection{Methodology Validation}

\textbf{DVRPC Methodology:} Industry-standard approach used by city planning departments for pedestrian infrastructure assessment

\textbf{Threshold Sensitivity:} Coverage thresholds (0.1, 0.4--0.8, 1.2) validated against visual inspection of sample roads

\textbf{Edge Detection Accuracy:} 5-foot offset buffers capture sidewalks along road edges while excluding distant sidewalks

\subsection{Output Validation}

\textbf{Data Integrity Checks:}
\begin{itemize}
    \item Sum of coverage categories equals total roads (3,245 + 620 + 521 = 4,386)
    \item TOD + Non-TOD roads sum to county total (1,117 + 3,269 = 4,386)
    \item All percentage calculations verified (e.g., 1,141 / 4,386 = 26.0\%)
    \item GeoJSON exports validated in QGIS (EPSG:4326 for web compatibility)
\end{itemize}

\section{Policy Recommendations}

\subsection{Immediate Priority Actions}

\begin{enumerate}
    \item \textbf{TOD No-Coverage Gap Closure:} Prioritize sidewalk installation on 502 TOD roads with no coverage to maximize transit connectivity
    \item \textbf{TOD One-Side Completion:} Complete 352 TOD roads with one-side coverage to achieve full both-sides pedestrian access
    \item \textbf{Unpaved Road Targeting:} Develop sidewalk installation strategy for unpaved roads (currently 4.0\% coverage)
\end{enumerate}

\subsection{Strategic Planning Recommendations}

\begin{itemize}
    \item \textbf{Equity Analysis:} Investigate socioeconomic factors associated with TOD vs non-TOD coverage gap
    \item \textbf{Transit Access Study:} Map pedestrian routes from residential areas to Metro-North stations
    \item \textbf{ADA Compliance Assessment:} Evaluate sidewalk quality and accessibility beyond binary coverage
    \item \textbf{Multi-Modal Integration:} Coordinate sidewalk expansion with bike lane and bus route planning
\end{itemize}

\subsection{Long-Term Vision}

\textbf{75\% TOD Coverage Target:} Achieving 75\% coverage in TOD areas would require completing 224 additional roads (from current 615 to 839 roads with coverage)

\textbf{50\% County-Wide Target:} Raising county-wide coverage to 50\% would require adding sidewalks to 1,052 roads (from 1,141 to 2,193 roads with coverage)

\textbf{Non-TOD Area Focus:} Closing the TOD vs non-TOD gap would require adding sidewalks to 1,178 non-TOD roads (to reach 54.9\% parity)

\section{Limitations and Future Enhancements}

\subsection{Data Limitations}

\begin{itemize}
    \item \textbf{Binary Coverage:} Analysis classifies coverage as present/absent but does not assess sidewalk quality, width, or ADA compliance
    \item \textbf{Temporal Snapshot:} County shapefiles represent a single point in time; does not track infrastructure change over time
    \item \textbf{Network Connectivity:} Does not analyze whether sidewalks form continuous pedestrian networks to key destinations
    \item \textbf{Socioeconomic Context:} Lacks demographic or economic data to assess equity implications
\end{itemize}

\subsection{Methodological Considerations}

\begin{itemize}
    \item \textbf{Threshold Sensitivity:} Coverage classification depends on ratio thresholds (0.1, 0.4--0.8, 1.2) - alternative thresholds may yield different results
    \item \textbf{Edge Buffer Size:} 5-foot offset buffers may miss sidewalks set back from road edges or capture unrelated infrastructure
    \item \textbf{TOD Buffer Distance:} 0.5-mile buffer is standard but actual walking catchment varies by topography, road crossings, and other barriers
\end{itemize}

\subsection{Future Enhancement Opportunities}

\begin{enumerate}
    \item \textbf{Network Analysis:} Assess pedestrian connectivity from residential areas to transit stations, schools, commercial centers
    \item \textbf{Quality Assessment:} Integrate sidewalk condition data (surface quality, width, ADA compliance, lighting)
    \item \textbf{Temporal Analysis:} Track sidewalk coverage changes over time using historical shapefiles
    \item \textbf{Equity Analysis:} Overlay demographic data to identify underserved populations
    \item \textbf{Multi-Modal Integration:} Combine sidewalk coverage with bike lane, bus route, and crosswalk data
    \item \textbf{Predictive Modeling:} Develop models to prioritize sidewalk investments based on demand, cost, and equity factors
\end{enumerate}

\section{Data and Code Availability}

\subsection{Machine-Readable Datasets (Excel)}

\textbf{Location:} \texttt{D:/Arcanum/Projects/Westchester/Output/Excel/}

\begin{itemize}
    \item \texttt{1\_EXECUTIVE\_SUMMARY.xlsx} - Answer front and center with complete breakdown (ONE SHEET)
    \item \texttt{2\_TOD\_COMPARISON.xlsx} - Side-by-side TOD vs Non-TOD analysis (ONE SHEET)
    \item \texttt{3\_ROAD\_TYPE\_ANALYSIS.xlsx} - Coverage by road type classification (ONE SHEET)
    \item \texttt{4\_AREA\_ANALYSIS.xlsx} - Coverage statistics by area in acres (ONE SHEET)
\end{itemize}

\textbf{Note:} All Excel files follow Druck standard with ONE SHEET per file.

\subsection{Geospatial Outputs (GeoJSON)}

\textbf{Location:} \texttt{D:/Arcanum/Projects/Westchester/Technical/data/processed/countywide\_sidewalk\_analysis/}

\textbf{County-Wide Files (EPSG:4326):}
\begin{itemize}
    \item \texttt{roads\_no\_coverage.geojson} - 3,245 roads without sidewalks
    \item \texttt{roads\_one\_side.geojson} - 620 roads with one-side coverage
    \item \texttt{roads\_both\_sides.geojson} - 521 roads with both-sides coverage
\end{itemize}

\textbf{TOD-Specific Files (EPSG:4326):}
\begin{itemize}
    \item \texttt{tod\_roads\_no\_coverage.geojson} - 502 TOD roads without sidewalks
    \item \texttt{tod\_roads\_one\_side.geojson} - 352 TOD roads with one-side coverage
    \item \texttt{tod\_roads\_both\_sides.geojson} - 263 TOD roads with both-sides coverage
    \item \texttt{metro\_north\_buffers.geojson} - 0.5-mile TOD buffer zones
\end{itemize}

\textbf{Total GeoJSON Size:} 596 MB (web-ready EPSG:4326 format for QGIS, ArcGIS, Leaflet, Mapbox)

\subsection{Statistical Summaries (JSON)}

\textbf{Location:} \texttt{D:/Arcanum/Projects/Westchester/Technical/data/processed/countywide\_sidewalk\_analysis/}

\begin{itemize}
    \item \texttt{county\_wide\_statistics.json} - Complete county-wide metrics including coverage counts, percentages, area statistics, and road type breakdowns
    \item \texttt{tod\_statistics.json} - TOD vs Non-TOD comparison metrics with detailed breakdowns
\end{itemize}

\subsection{Analysis Code}

\textbf{Location:} \texttt{D:/Arcanum/Projects/Westchester/Technical/src/}

\textbf{Core Analysis Scripts:}
\begin{itemize}
    \item \texttt{countywide\_sidewalk\_analyzer.py} - Main analysis engine (652 lines) implementing DVRPC methodology
    \item \texttt{transit\_sidewalk\_analyzer.py} - TOD buffer analysis and TOD vs non-TOD comparison
    \item \texttt{generate\_excel\_reports.py} - Druck-compliant Excel report generator
\end{itemize}

\textbf{Reproduction Instructions:} See \texttt{Technical/README.md} for step-by-step execution guide

\section{Conclusion}

\textbf{Primary Answer:} 54.9\% of roads within 0.5 miles of Metro-North stations have sidewalk coverage (615 out of 1,117 roads), representing MODERATE ADEQUACY with substantial room for improvement.

\textbf{Key Insights:}
\begin{itemize}
    \item TOD areas have 3.0x better coverage than non-TOD areas (54.9\% vs 18.2\%)
    \item 45.1\% of TOD roads (502 roads) still lack sidewalks - priority gap for transit access
    \item 83.9\% of non-TOD roads lack coverage - severe pedestrian infrastructure inequality
    \item Unpaved roads show only 4.0\% coverage - systematically underserved road type
\end{itemize}

\textbf{Research Contribution:} This analysis provides the first comprehensive spatial assessment of sidewalk coverage in Westchester County using industry-standard DVRPC methodology with 1-side vs 2-side detection.

\textbf{Practical Value:} Results enable data-driven prioritization of sidewalk investments to maximize transit connectivity, pedestrian safety, and infrastructure equity.

\end{document}
